% Shared titles: update each definition to change both the TOC and first slide title.
\newcommand{\lossestitle}{Fury's Attack}
\newcommand{\sakaartitle}{The Winter Soldier}
\newcommand{\decisionstitle}{The Choice to Resist}
\newcommand{\prioritiestitle}{A Promise to Bucky}
\newcommand{\returntitle}{The Fall of S.H.I.E.L.D.}
\newcommand{\modeltitle}{A Tally of Those Protected}

{
    \setbeamertemplate{footline}{}
    \begin{frame}[plain,noframenumbering]
        \titlepage
        \note{Introduce Steve's journey from Fury's attack to his search for Bucky.}
    \end{frame}
}

\begin{frame}{\contentslabel}{From a hidden conspiracy to an unfinished reunion}
    \begin{SlideContents}
    \end{SlideContents}
    \note{Preview the conspiracy, HYDRA's defeat, and Steve's search for Bucky.}
\end{frame}

\section{The Conspiracy}
\subsection{\lossestitle}
\begin{frame}{\lossestitle}{Fury's attack and a threat within S.H.I.E.L.D.}
    % The image is scaled to its specified height and cropped to the column width.
    \begin{SlideBulletsPhoto}{assets/winter_soldier_street.jpg}{4.2cm}{The Winter Soldier on the street.}
        \begin{itemize}
            \item \textbf{Nick Fury} is attacked by a masked assassin.
            \item \textbf{S.H.I.E.L.D.} brands Steve a fugitive.
            \item \textbf{Steve} discovers a threat within the agency.
        \end{itemize}
    \end{SlideBulletsPhoto}
    \note{Fury's attack draws Steve into a conspiracy within S.H.I.E.L.D.}
\end{frame}

\subsection{\sakaartitle}
\begin{frame}{\sakaartitle}{The assassin is someone Steve once knew}
    \begin{SlideTitlePhoto}{assets/captain_vs_winter_soldier.jpg}{5.2cm}{0.95\linewidth}
        % Scale to the specified height and crop any horizontal overflow.
        \caption{Captain America faces the Winter Soldier.}
    \end{SlideTitlePhoto}
    \note{The Winter Soldier is Bucky Barnes, Steve's friend under HYDRA's control.}
\end{frame}

\begin{frame}{A Friend Behind the Mask}{Steve and Natasha uncover HYDRA's infiltration}
    \begin{SlidePhotoPair}
        % Both paired images use the same height and width limits.
        \SlidePhotoPanel{assets/captain_and_black_widow.jpg}{3.8cm}{Steve and Natasha search for the truth.}
        \hfill
        \SlidePhotoPanel{assets/shield_helicarriers.jpg}{3.8cm}{Project Insight threatens those HYDRA considers enemies.}
        \caption{The allies uncover a conspiracy within S.H.I.E.L.D.}
    \end{SlidePhotoPair}
    \note{Steve, Natasha, and Sam follow the evidence to HYDRA.}
\end{frame}

\subsubsection{The Friend Behind the Mask}
\begin{frame}{The Friend Behind the Mask}{Bucky has been turned into HYDRA's weapon}
    \begin{SlideTitlePhoto}{assets/captain_and_bucky_by_river.jpg}{5.2cm}{0.95\linewidth}
        % Scale to the specified height and crop any horizontal overflow.
        \caption{Steve refuses to give up on Bucky.}
    \end{SlideTitlePhoto}
    \note{Steve recognizes his friend beneath HYDRA's conditioning.}
\end{frame}

\subsection{\decisionstitle}
\begin{frame}{\decisionstitle}{Project Insight turns security into a weapon}
    \begin{SlideThreeColumns}
        \begin{column}{0.30\textwidth}
            \textbf{What Steve Learns}
            \begin{itemize}
                \item HYDRA is inside S.H.I.E.L.D.
                \item Bucky is the Winter Soldier
                \item Project Insight targets threats
            \end{itemize}
        \end{column}
        \begin{column}{0.30\textwidth}
            \textbf{The Plan}
            \begin{enumerate}
                \item Expose HYDRA
                \item Stop the helicarriers
                \item Protect the people
            \end{enumerate}
        \end{column}
        \begin{column}{0.30\textwidth}
            \textbf{The Choice}
            \begin{description}
                \item[Steve] \keyphrase{Protect the people}
                \item[Natasha] Reveal the files
                \item[Sam] Help stop the carriers
            \end{description}
        \end{column}
    \end{SlideThreeColumns}
    \note{Steve chooses to stop Project Insight and oppose the agency he serves.}
\end{frame}

\begin{frame}{The Road to the Triskelion}{From Fury's attack to the helicarrier launch}
    \SlideTimeline
    \begin{tikzpicture}[
            stage/.style={draw=offwhite,align=center,minimum width=3cm,minimum height=1cm},
            every node/.style={font=\small}]
        \node[stage] (loss) {Attack\\Fury};
        \node[stage,right of=loss,node distance=4cm] (exile) {Truth\\HYDRA};
        \node[stage,final stage,right of=exile,node distance=4cm] (return) {Triskelion\\Stop Insight};
        \draw[progress arrow] (loss) -- (exile);
        \draw[progress arrow] (exile) -- (return);
    \end{tikzpicture}
    \note{Trace Steve's journey from Fury's attack to the Triskelion.}
\end{frame}

\section{Steve's Choice}
{
    \setbeamertemplate{frametitle}{}
    \begin{frame}
        \begin{tikzpicture}[remember picture,overlay]
            \node[inner sep=0pt,font=\LARGE] at (current page.center) {A friend is worth finding.};
        \end{tikzpicture}
        \note{Steve refuses to abandon Bucky, even after HYDRA's conditioning.}
    \end{frame}
}

\subsection{\prioritiestitle}
\begin{frame}{\prioritiestitle}
    \centering
    {\Large Bucky may have forgotten their past. \keyphrase{Steve has not.}\par}
    \medskip
    {\footnotesize From the historical records: \citet{lore_thor,lore_loki,lore_asgard}.\par}
    \note{Emphasize Steve's loyalty to Bucky.}
\end{frame}

\begin{frame}{A Promise Kept}{Steve will not leave Bucky behind}
    \begin{SlideTable}
        \caption{The path from HYDRA's control to an uncertain reunion.}
        \setlength\extrarowheight{4pt}
        \begin{tabularx}{\linewidth}{lYY}
            \toprule
            \rowcolor{softblack}
            \textbf{Stage} & \textbf{Disruption} & \textbf{Response}\\
            \midrule
            S.H.I.E.L.D. & HYDRA hides inside the agency & Follow the evidence\\
            Triskelion & Project Insight is about to launch & Stop the helicarriers\\
            River & Bucky fights Steve, then saves him & Hold on to their friendship\\
            \midrule
            \highlightrow
            Resolution & Bucky disappears with his past unresolved & Steve sets out to find him\\
            \bottomrule
        \end{tabularx}
    \end{SlideTable}
    \note{Follow the crisis from HYDRA's infiltration to Steve's search for Bucky.}
\end{frame}

\subsection{\returntitle}
\begin{frame}{\returntitle}{HYDRA is exposed, and the agency collapses}
    \begin{block}{The Threat}
        HYDRA plans to eliminate people it considers dangerous.
    \end{block}
    \begin{alertblock}{The Resistance}
        Steve, Natasha, and Sam expose HYDRA and stop the helicarriers.
    \end{alertblock}
    \begin{exampleblock}{The Aftermath}
        S.H.I.E.L.D. falls. Bucky saves Steve, then disappears.\footnote{A chronicle of the conspiracy and its aftermath.}
    \end{exampleblock}
    \medskip
    From the historical records: \citet{lore_thor}.
    \note{HYDRA's plan is stopped, but Bucky's future remains unresolved.}
\end{frame}

\subsection{\modeltitle}
\begin{frame}[fragile]{\modeltitle}{An imagined count of those protected}
    \begin{SlideCodeExplanation}
        \begin{column}{0.46\textwidth}
            \begin{Verbatim}[fontsize=\footnotesize]
def protected_fraction(p, N):
    if N <= 0 or p < 0 or p > N:
        raise ValueError("Bad input")
    return p / N

saved, total = 75, 100
fraction = protected_fraction(saved, total)
print(f"Saved: {fraction:.0%}")
      \end{Verbatim}
        \end{column}
        \begin{column}{0.46\textwidth}
            \textbf{This imagined tally} represents people protected from Project Insight.

            \medskip
            This is an imagined count, not a number given in the story:
            \[
                r=\frac{p}{N},\qquad N>0,\quad 0\leq p\leq N.
            \]
        \end{column}
    \end{SlideCodeExplanation}
    \note{The 75 percent figure is an imagined example, not a count from the story.}
\end{frame}

\begin{frame}{A Narrowing Window to Act}{A hypothetical account of stopping Project Insight}
    This imagined model estimates the chance of stopping Project Insight:
    \begin{equation}
        P(E)=\frac{1}{1+\exp[-(\alpha C-\gamma D)]},\qquad \alpha,\gamma>0.
    \end{equation}
    Here, $C$ is the available response capacity and $D$ is the delay before the response begins.
    This model is imagined; the chronicle gives no exact response data.

    \medskip
    Greater response capacity and an earlier intervention improve the chance of stopping the helicarriers.
    \note{More capacity and less delay improve the chance of success in this imagined model.}
\end{frame}

\section{\referenceslabel}
\setbeamertemplate{frametitle continuation}[from second]
\begin{frame}[allowframebreaks]{\referenceslabel}{Records of HYDRA and the Winter Soldier}
    \begin{SlideReferences}
    \bibliographystyle{plainnat}
    \bibliography{references}
    \end{SlideReferences}
    \note{Point to the fictional sources listed here for the chronicle.}
\end{frame}
